\chapter{Tree Decompositions}\label{ch:treedecomp}

The techniques of the previous chapter answer questions about one root path, one subtree or one fixed root. Harder questions
need the tree cut into pieces with a logarithmic guarantee. Paths between arbitrary vertices under updates call for heavy-light
decomposition; problems about \emph{all} paths of a tree at once call for centroid decomposition; statistics of every subtree
can be gathered in one shared table with DSU on tree, or in linear time with long-path decomposition; many small queries shrink
to virtual trees; and trees that change shape are kept as Euler tour trees or link-cut trees. Most figures use the same running
example, so the decompositions can be compared directly.

\begin{chapterintro}{\plat\ \ioi}{Chapter~\ref{ch:trees} (Euler tours, LCA, small-to-large) and Chapter~\ref{ch:segtree}
(segment trees).}
  \item split a tree into heavy paths and answer path queries with updates in $\Oh{\log^2 n}$;
  \item decompose a tree by centroids and solve problems about all of its paths, offline and online;
  \item gather subtree statistics in one shared table (DSU on tree), and depth-indexed DP in linear time (long-path
        decomposition);
  \item compress a tree to a few key vertices (virtual trees);
  \item keep a forest under link and cut (Euler tour trees), and know what link-cut trees and top trees add.
\end{chapterintro}

% the running example: 13 vertices rooted at 0
\newcommand{\rexcoords}{%
  \foreach \nm/\px/\py in {0/3.0/0, 1/1.2/-1.1, 2/4.4/-1.1, 3/6.0/-1.1, 4/0.6/-2.2, 5/1.8/-2.2, 6/4.4/-2.2, 7/0.0/-3.3,
      8/1.2/-3.3, 9/3.8/-3.3, 10/5.0/-3.3, 11/0.0/-4.4, 12/5.0/-4.4}
    \coordinate (R\nm) at (\px, \py);}
\newcommand{\rexedges}{(R0) -- (R1) -- (R4) -- (R7) -- (R11) (R2) -- (R6) -- (R10) -- (R12) (R0) -- (R2) (R0) -- (R3)
  (R1) -- (R5) (R4) -- (R8) (R6) -- (R9)}

% ======================================================================================================================
\section[Heavy-light decomposition]{Heavy-light decomposition \hfill \plat}\label{sec:hld}

Consider Path Queries~II: vertex values change, and a query asks for the maximum on the path between two arbitrary vertices. The
Euler tour of Section~\ref{sec:eulertour} handles subtrees and sums along root paths, but a maximum cannot be ``subtracted'', and a
path between two arbitrary vertices is not a prefix of anything. \defn{Heavy-light decomposition} (HLD) cuts the tree into
vertical paths such that every path of the tree uses only a few of them, and lays each vertical path out contiguously in an
array. A segment tree over that array then answers each piece with one range query.

\paragraph{Heavy and light edges.} For every vertex with children, the \defn{heavy child} is a child with the largest subtree
(ties broken arbitrarily); the edge to it is \emph{heavy}, and the edges to all other children are \emph{light}. Every vertex has at
most one heavy child edge, so the heavy edges form disjoint vertical paths, the \defn{heavy paths} or \emph{chains}. Every vertex
lies on exactly one chain, and the top vertex of its chain is its \emph{head}; a vertex whose parent edge is light starts a new
chain (Figure~\ref{fig:hld}).

\begin{figure}[htbp]
\centering
\begin{tikzpicture}
  \tikzset{pl/.style={font=\sffamily\scriptsize\bfseries, text=fgblue, inner sep=1.5pt}}
  \rexcoords
  \draw[hed] (R0) -- (R1) -- (R4) -- (R7) -- (R11);
  \draw[bed] (R2) -- (R6) -- (R10) -- (R12);
  \draw[xed] (R0) -- (R2) (R0) -- (R3) (R1) -- (R5) (R4) -- (R8) (R6) -- (R9);
  \foreach \nm/\sy/\ps/\an in {0/hot/0/above, 1/hot/1/left, 4/hot/2/left, 7/hot/3/left, 11/hot/4/left, 8/n/5/right,
      5/n/6/right, 3/n/7/right, 2/cool/8/right, 6/cool/9/right, 10/cool/10/right, 12/cool/11/right, 9/n/12/left}
    \node[vtx, \sy, label={[pl]\an:\ps}] (V\nm) at (R\nm) {\nm};
  \draw[hed] (7.0, -0.55) -- ++(0.55, 0); \node[lbl, anchor=west] at (7.65, -0.55) {heavy edge};
  \draw[xed] (7.0, -0.95) -- ++(0.55, 0); \node[lbl, anchor=west] at (7.65, -0.95) {light edge};
  \node[pl] at (7.27, -1.35) {7}; \node[lbl, anchor=west] at (7.65, -1.35) {position of the vertex};
  \legendat{(7.27, -1.8)}{hot/chain with head 0, cool/chain with head 2, n/one-vertex chains}
  \node[blbl, anchor=east] at (-0.2, -5.75) {vertex};
  \arrs{H}{0}{-5.75}{0/hot, 1/hot, 4/hot, 7/hot, 11/hot, 8/n, 5/n, 3/n, 2/cool, 6/cool, 10/cool, 12/cool, 9/n}
  \idx{H}{13}
  \node[lbl, anchor=east] at (-0.2, -6.28) {position};
  \draw[mbrace] ($(H-2.south west)+(0, -0.45)$) -- node[below=6pt, lbl] {subtree of 1: positions $[1, 7)$}
    ($(H-7.south east)+(0, -0.45)$);
  \draw[mbrace] ($(H-9.south west)+(0, -0.45)$) -- node[below=6pt, lbl] {subtree of 2: $[8, 13)$}
    ($(H-13.south east)+(0, -0.45)$);
\end{tikzpicture}
\caption{Heavy-light decomposition of the running example. Heavy edges form the chains 0--1--4--7--11 and 2--6--10--12; every
other vertex is a chain by itself. Numbering each chain from its head, and only then the light subtrees hanging from it, deepest
first, makes every chain \emph{and} every subtree a contiguous block of positions.}
\label{fig:hld}
\end{figure}

\begin{fact}
On the path from any vertex to the root there are at most $\log_2 n$ light edges. If $p \to v$ is a light edge, the heavy child
$h$ of $p$ has $\mathit{sub}(h) \ge \mathit{sub}(v)$, so $\mathit{sub}(p) \ge 1 + \mathit{sub}(v) + \mathit{sub}(h) > 2\,
\mathit{sub}(v)$: every light edge on the way up more than doubles the subtree size, which never exceeds $n$
(Figure~\ref{fig:hldlight}).
\end{fact}

\begin{figure}[htbp]
\centering
\begin{tikzpicture}
  \node[vtx] (p) at (0, 0) {$p$};
  \node[vtx, cool] (v) at (-1.5, -1.2) {$v$};
  \node[vtx, hot] (h) at (1.5, -1.2) {$h$};
  \draw[ed, dashed] (p) -- ++(0, 0.8);
  \draw[xed] (p) -- (v);
  \draw[hed] (p) -- (h);
  \draw[draw=fgblue, fill=fgsky!18, line width=0.6pt] (v.south) -- ++(-0.55, -1.0) -- ++(1.1, 0) -- cycle;
  \draw[draw=fgverm, fill=fgverm!10, line width=0.6pt] (h.south) -- ++(-0.8, -1.5) -- ++(1.6, 0) -- cycle;
  \node[font=\sffamily\scriptsize, text=fgblue] at (-1.5, -2.85) {$\mathit{sub}(v) = s$};
  \node[font=\sffamily\scriptsize, text=fgverm] at (1.5, -3.35) {$\mathit{sub}(h) \ge s$};
  \node[font=\sffamily\scriptsize, anchor=west] at (0.3, 0.3) {$\mathit{sub}(p) \ge 2s + 1$};
  \begin{scope}[xshift=5.4cm]
    \foreach \nm/\py in {r/0.3, u4/-0.45, u3/-1.2, u2/-1.95, u1/-2.7, x/-3.45}
      \node[vtxs] (\nm) at (0, \py) {};
    \node[font=\sffamily\scriptsize] at (x) {$x$};
    \draw[hed] (r) -- (u4) (u2) -- (u1);
    \draw[xed] (u4) -- (u3) (u3) -- (u2) (u1) -- (x);
    \node[lbl, anchor=west] at (x.east) {$\mathit{sub} = 1$};
    \node[lbl, anchor=west] at (u1.east) {$\ge 3$};
    \node[lbl, anchor=west] at (u3.east) {$\ge 7$};
    \node[lbl, anchor=west] at (u4.east) {$\ge 15$};
    \node[lbl, anchor=west] at (r.east) {root};
    \node[notel, anchor=west, text width=4.3cm] at (1.5, -1.6) {Climbing from $x$, every light edge (dashed) more than
      doubles the subtree size; heavy edges never shrink it. A root path therefore has at most $\log_2 n$ light edges.};
  \end{scope}
\end{tikzpicture}
\caption{The light-edge lemma behind every $\log n$ bound of this section.}
\label{fig:hldlight}
\end{figure}

So a root path touches at most $\log_2 n + 1$ chains, and the path between $u$ and $v$, which lies inside the union of the two
root paths, is cut into at most $2\log_2 n + 2$ \emph{chain pieces}.

\paragraph{Positions.} Number the vertices so that every chain is a contiguous block that starts at its head: walk a chain from
its head, numbering as you go, and push the light children met on the way onto a stack, to be numbered later. Because the light
subtrees are taken from the stack deepest first, the numbering is also a DFS preorder, so the subtree of $v$ is the block
$[\mathit{pos}(v), \mathit{pos}(v) + \mathit{sub}(v))$ exactly as with the Euler tour (Figure~\ref{fig:hld}). One array, one
segment tree, serves both path queries and subtree queries.

\paragraph{Answering a path query.} While $a$ and $b$ lie on different chains, look at the one whose chain head is
\emph{deeper}, say $b$: its chain piece from $\mathit{head}(b)$ down to $b$ is the position range $[\mathit{pos}(\mathit{head}(b)),
\mathit{pos}(b)]$, so query it and move $b$ to the parent of its head. Once both are on one chain, the piece between them is a
single range (Figure~\ref{fig:hldpath}). Moving the vertex with the deeper head never passes the LCA: if the LCA were on $b$'s
chain at or below $\mathit{head}(b)$, then $a$, a descendant of it on another chain, would have a head strictly deeper than
$\mathit{head}(b)$. Every move crosses a light edge, so there are $\Oh{\log n}$ pieces, and with a segment tree a query costs
$\Oh{\log^2 n}$ and an update $\Oh{\log n}$.

\begin{figure}[htbp]
\centering
\begin{tikzpicture}
  \tikzset{pc/.style={circle, draw=fgverm, fill=white, line width=0.8pt, inner sep=1pt, minimum size=4.2mm,
    font=\sffamily\scriptsize\bfseries, text=fgverm}}
  \rexcoords
  \draw[line width=5pt, draw=fggreen, opacity=0.3] (R11) -- (R7) -- (R4) -- (R1) -- (R0) -- (R2) -- (R6) -- (R9);
  \draw[ed] \rexedges;
  \foreach \nm/\sy in {0/n, 1/n, 2/n, 3/n, 4/n, 5/n, 6/n, 7/n, 8/n, 9/good, 10/n, 11/good, 12/n}
    \node[vtx, \sy] at (R\nm) {\nm};
  \node[pc] at ($(R9)+(-0.6, -0.3)$) {1};
  \node[pc] at ($(R2)!0.5!(R6)+(-0.42, 0)$) {2};
  \node[pc] at ($(R1)!0.5!(R4)+(-0.5, 0.05)$) {3};
  \node[notel, anchor=north west, text width=5.6cm] at (6.9, 0.2) {\textbf{The query 11--9.}\\[3pt]
    \textcolor{fgverm}{\textbf{1}}\quad The heads are 0 and 9; the head of 9 is deeper, so take the range $[12, 13)$ and move
    to $\mathit{parent}(9) = 6$.\\[2pt]
    \textcolor{fgverm}{\textbf{2}}\quad The heads are 0 and 2; the head of 6 is deeper, so take $[8, 10)$ and move to
    $\mathit{parent}(2) = 0$.\\[2pt]
    \textcolor{fgverm}{\textbf{3}}\quad Now 0 and 11 share a chain: take $[0, 5)$.};
  \node[blbl, anchor=east] at (-0.2, -5.75) {vertex};
  \arrs{H}{0}{-5.75}{0/good, 1/good, 4/good, 7/good, 11/good, 8/n, 5/n, 3/n, 2/good, 6/good, 10/n, 12/n, 9/good}
  \idx{H}{13}
  \cellframe{H-1}{H-5}{fgverm}
  \cellframe{H-9}{H-10}{fgverm}
  \cellframe{H-13}{H-13}{fgverm}
  \node[pc] at ($(H-3.north)+(0, 0.4)$) {3};
  \node[pc] at ($(H-9.north east)+(0, 0.4)$) {2};
  \node[pc] at ($(H-13.north)+(0, 0.4)$) {1};
\end{tikzpicture}
\caption{A path query on the running example. The path from 11 to 9 (green) is cut into three chain pieces, which are three
ranges of positions (framed); each costs one segment-tree query.}
\label{fig:hldpath}
\end{figure}

\begin{problem}{Path Queries II (CSES)}
A tree has a value at every vertex. Process $q$ operations: set the value of a vertex, or report the maximum value on the path
between two vertices ($n, q \le 2 \cdot 10^5$).
\solution HLD positions, and a bottom-up max segment tree over them. The program below runs in 0.25\,s on a path and 0.21\,s on a
random tree with $n = q = 2 \cdot 10^5$ (\cc{g++ -O2}).
\end{problem}

\begin{code}
void solve() {  // HLD (cphld) is pasted above solve()
    int n, q;
    cin >> n >> q;
    vi val(n);
    for (ll& x : val) cin >> x;
    vvi adj(n);
    rep(i, n - 1) {
        int a, b;
        cin >> a >> b;
        a--, b--;
        adj[a].push_back(b);
        adj[b].push_back(a);
    }
    HLD h(adj);
    // a bottom-up max segment tree over the HLD positions
    vi t(2 * n);
    rep(v, n) t[n + h.pos[v]] = val[v];
    per(i, 1, n) t[i] = max(t[2 * i], t[2 * i + 1]);
    auto query = [&](int l, int r) {  // max over positions [l, r)
        ll res = LLONG_MIN;
        for (l += n, r += n; l < r; l >>= 1, r >>= 1) {
            if (l & 1) chmax(res, t[l++]);
            if (r & 1) chmax(res, t[--r]);
        }
        return res;
    };
    while (q--) {
        int type, a, b;
        cin >> type >> a >> b;
        if (type == 1) {  // the value of vertex a becomes b
            int i = n + h.pos[a - 1];
            for (t[i] = b; i > 1; i >>= 1) t[i >> 1] = max(t[i], t[i ^ 1]);
        } else {  // O(log n) chain pieces, one range query each
            ll best = LLONG_MIN;
            h.for_path(a - 1, b - 1,
                       [&](int l, int r) { chmax(best, query(l, r)); });
            print(best);
        }
    }
}
\end{code}

\usetemplate{Heavy-light decomposition}{Ctrl+Alt+Y, H}{cphld}

The template's \cc{for_path(a, b, f)} calls \cc{f(l, r)} once for every piece, as a half-open range of positions, so any
segment tree plugs in: sums, minima, a lazy tree with range assignment. It also offers \cc{lca(a, b)}, which runs the same loop
without the queries and is one of the fastest LCA methods in practice: $\Oh{\log n}$ per query, $\Oh{n}$ memory, and on most
trees only a few iterations, one per chain switch.

\begin{pitfall}[Values on edges]
When the values live on edges, store each edge's value at its deeper endpoint. A path query must then skip the LCA, whose stored
value belongs to the edge above it: in the last piece, use $[\mathit{pos}(\mathrm{lca}) + 1, \mathit{pos}(v)]$. The template marks
the line to change. Forgetting it gives answers that are wrong only when the edge above the LCA is the extreme one, which random
tests rarely catch.
\end{pitfall}

\begin{trick}[Subtrees and paths in one structure]
Because HLD positions are also a preorder, ``add $x$ to every value in the subtree of $v$'' is a range update on
$[\mathit{pos}(v), \mathit{pos}(v) + \mathit{sub}(v))$, while ``sum on the path $u$--$v$'' is the usual $\Oh{\log n}$ ranges. One
lazy segment tree handles both kinds of operations, which neither the Euler tour alone nor binary lifting can do.
\end{trick}

\begin{recognize}
``Update a vertex or edge; query a sum, minimum or maximum on the path between two vertices'', or range updates on paths. For
path queries without updates, binary lifting or prefix sums along the Euler tour are simpler; HLD earns its place where paths
\emph{and} updates meet. With $n = q = 2 \cdot 10^5$ the $\Oh{\log^2 n}$ bound is far from tight: most paths cross only a few
light edges.
\end{recognize}

% ======================================================================================================================
\section[Centroid decomposition]{Centroid decomposition \hfill \plat}\label{sec:centroid}

Some problems ask about \emph{all} paths of a tree: how many have exactly $k$ edges, how many have weight at most $W$, which has
the largest XOR. There are $\Theta(n^2)$ paths, but a divide-and-conquer handles them in bulk: pick a vertex $c$, deal with all
paths through $c$, delete $c$, and recurse into the pieces that remain. Every path is handled exactly once, at the first of its
vertices to be picked. Picking a \emph{centroid} each time keeps the recursion $\Oh{\log n}$ deep.

\paragraph{Centroids.} A \defn{centroid} of a tree with $n$ vertices is a vertex whose removal leaves components of at most $n/2$
vertices each. One always exists, and a walk finds it: root the tree anywhere and, while some child of the current vertex has a
subtree of more than $n/2$ vertices, step into that child. After a step into $c$, the component above $c$ has $n -
\mathit{sub}(c) < n/2$ vertices, so only children can be too big; subtree sizes decrease, so the walk stops, at a centroid. A tree
has one centroid or two adjacent ones.

\paragraph{The decomposition.} Find a centroid $c$, process the paths through $c$, delete $c$, and recurse into every remaining
component. Each component has at most half the vertices of the one it came from, so the recursion is at most $\lfloor \log_2 n
\rfloor + 1$ levels deep, and each vertex belongs to one component per level. If the work at each component is linear in its
size, the total is $\Oh{n \log n}$. Linking every centroid below the centroid whose deletion created its component gives the
\defn{centroid tree} (Figure~\ref{fig:centroid}).

\begin{figure}[htbp]
\centering
\begin{minipage}[c]{6.9cm}\centering
\begin{tikzpicture}
  \rexcoords
  \draw[ed] \rexedges;
  \foreach \nm/\sy in {0/hot, 3/warn, 4/warn, 6/warn, 1/cool, 2/cool, 7/cool, 8/cool, 9/cool, 10/cool, 5/n, 11/n, 12/n}
    \node[vtx, \sy] at (R\nm) {\nm};
  \foreach \sy/\tx [count=\ci] in {hot/level 0, warn/level 1, cool/level 2, n/level 3} {
    \node[draw, \sy, minimum size=2.6mm, inner sep=0pt] (lg\ci) at ({1.6*(\ci - 1)}, -5.25) {};
    \node[lbl, anchor=west] at ($(lg\ci.east)+(0.08, 0)$) {\tx};
  }
\end{tikzpicture}
\end{minipage}
\hspace{4mm}
\begin{minipage}[c]{6cm}\centering
\begin{forest}
  cptree, for tree={s sep=3mm, l sep=7mm}
  [0, hot [4, warn [1, cool [5]] [7, cool [11]] [8, cool]] [6, warn [2, cool] [9, cool] [10, cool [12]]] [3, warn]]
\end{forest}
\end{minipage}
\caption{Centroid decomposition of the running example. Removing the centroid 0 leaves parts of 6, 5 and 1 vertices, at most
$13/2$ each; their centroids 4, 6 and 3 form level 1, and so on. Right: the centroid tree, of depth $3 \le \log_2 13$. The path
from 11 to 9 contains exactly one vertex of the lowest level, 0, which is the LCA of 11 and 9 in the centroid tree.}
\label{fig:centroid}
\end{figure}

\begin{idea}
For any two vertices $u$ and $v$, the path between them contains exactly one vertex of minimum level: their LCA $c$ in the
centroid tree. Until $c$ is deleted, $u$ and $v$ stay in one component, and $c$ is the first vertex of the path to be deleted.
Hence the ``paths through $c$ inside $c$'s component'' partition all paths, and $\mathrm{dist}(u, v) = \mathrm{dist}(u, c) +
\mathrm{dist}(c, v)$.
\end{idea}

\begin{problem}{Fixed-Length Paths I (CSES)}
Count the paths with exactly $k$ edges in a tree of $n \le 2 \cdot 10^5$ vertices.
\solution At a centroid $c$, a path through $c$ joins two vertices in different branches of $c$ (or $c$ itself) whose depths below
$c$ add up to $k$. Keep $\mathit{cnt}[d]$, the number of vertices at depth $d$ found so far, starting with $c$ itself at depth
0. Visit the branches one by one: first add $\mathit{cnt}[k - d]$ for every depth $d$ of the branch, which pairs it only with
earlier branches, then add the branch's own depths to $\mathit{cnt}$ (Figure~\ref{fig:fixedlen}).
\end{problem}

\begin{figure}[htbp]
\centering
\begin{minipage}[c]{4.1cm}\centering
\begin{tikzpicture}
  \node[vtx, hot] (c) at (1.4, 0) {$c$};
  \foreach \nm/\px/\py/\sy/\dv in {a/0/-1.0/cool/1, a1/-0.5/-2.0/cool/2, a2/0.5/-2.0/cool/2, b/1.4/-1.0/warn/1,
      b1/1.4/-2.0/warn/2, d/2.8/-1.0/pur/1, d1/2.8/-2.0/pur/2, d2/2.8/-3.0/pur/3}
    \node[vtx, \sy] (\nm) at (\px, \py) {\dv};
  \draw[ed] (c) -- (a) (a) -- (a1) (a) -- (a2) (c) -- (b) (b) -- (b1) (c) -- (d) (d) -- (d1) (d1) -- (d2);
  \foreach \tx/\px in {A/0, B/1.4, C/2.8} \node[blbl] at (\px, -3.65) {\tx};
\end{tikzpicture}
\end{minipage}
\hspace{4mm}
\begin{minipage}[c]{10.2cm}\centering\small
\begin{tabular}{@{}lccccr@{}}
\toprule
 & $\mathit{cnt}[0]$ & $\mathit{cnt}[1]$ & $\mathit{cnt}[2]$ & $\mathit{cnt}[3]$ & new paths \\
\midrule
start: $c$ alone & 1 & 0 & 0 & 0 & \\
after branch A & 1 & 1 & 2 & 0 & 0 \\
after branch B & 1 & 2 & 3 & 0 & $2 + 1 = 3$ \\
after branch C & 1 & 3 & 4 & 1 & $3 + 2 + 1 = 6$ \\
\bottomrule
\end{tabular}
\end{minipage}
\caption{Counting the paths of length $k = 3$ through a centroid $c$; vertices show their depth below $c$. Branch B finds
$\mathit{cnt}[2] + \mathit{cnt}[1] = 3$ partners for its depths 1 and 2, branch C finds 6, so $0 + 3 + 6 = 9$ paths pass through
$c$. Paths inside one branch are counted later, at a deeper centroid.}
\label{fig:fixedlen}
\end{figure}

\begin{code}
void solve() {
    int n, k;
    cin >> n >> k;
    vvi adj(n + 1);
    rep(i, n - 1) {
        int a, b;
        cin >> a >> b;
        adj[a].push_back(b);
        adj[b].push_back(a);
    }
    vector<bool> gone(n + 1, false);  // centroids already removed
    vi par(n + 1), sub(n + 1), dep(n + 1), cnt(n + 1, 0);
    auto bfs = [&](int s, int p) {  // the part still attached to s, BFS order
        vi order{s};
        par[s] = p;
        rep(i, sz(order)) for (int v : adj[order[i]]) {
            if (v == par[order[i]] || gone[v]) continue;
            par[v] = order[i];
            order.push_back(v);
        }
        return order;
    };
    ll ans = 0;
    vi roots{1};  // one vertex of every component still to decompose
    while (!roots.empty()) {
        int r = roots.back();
        roots.pop_back();
        vi order = bfs(r, 0);
        per(i, sz(order)) {
            int v = order[i];
            sub[v] = 1;
            for (int u : adj[v])
                if (u != par[v] && !gone[u]) sub[v] += sub[u];
        }
        int total = sz(order), c = r;
        for (bool moved = true; moved;) {  // step into a child with > half
            moved = false;
            for (int u : adj[c])
                if (u != par[c] && !gone[u] && 2 * sub[u] > total) {
                    c = u, moved = true;
                    break;
                }
        }
        // paths through c: pair each branch with c and the earlier branches
        cnt[0] = 1, dep[c] = 0;
        int maxd = 0;
        for (int u : adj[c]) {
            if (gone[u]) continue;
            vi part = bfs(u, c);
            for (int v : part) dep[v] = dep[par[v]] + 1;
            for (int v : part)
                if (dep[v] <= k) ans += cnt[k - dep[v]];
            for (int v : part)
                if (dep[v] <= k) cnt[dep[v]]++, chmax(maxd, dep[v]);
        }
        fill(cnt.begin(), cnt.begin() + maxd + 1, 0);
        gone[c] = true;
        for (int u : adj[c])
            if (!gone[u]) roots.push_back(u);
    }
    print(ans);
}
\end{code}

The decomposition uses an explicit stack of component roots and BFS orders, so it never recurses: on a path of $2 \cdot 10^5$
vertices, the worst case for depth, it runs in 0.20\,s, and in 0.17\,s on a random tree. The counting array is cleared only up to
the largest depth used, so each component costs time linear in its size, never $\Oh{n}$.

\paragraph{Fixed-Length Paths II: a range of lengths.} When the length may be anywhere in $[k_1, k_2]$, the lookup
$\mathit{cnt}[k - d]$ becomes a range sum over $\mathit{cnt}[k_1 - d .. k_2 - d]$, so the counts need prefix sums. Rebuilding the
prefix sums after every branch costs the current largest depth, which could be large while the branch is tiny; the fix is to
visit the branches in increasing order of height. Then the largest depth so far never exceeds the height of the branch being
added, so each rebuild costs no more than that branch's size, and the whole decomposition stays $\Oh{n \log n}$ (the sort by
height costs $\Oh{\deg \log \deg}$ per centroid).

\begin{trick}[Inclusion--exclusion instead of branch order]
For conditions that a sort handles, such as ``depth sum at most $K$'', count \emph{all} pairs in the component, ignoring branches,
with a sort and two pointers; the pairs inside one branch were counted wrongly (their real path does not pass through $c$), so
subtract the same count run on each branch's list of depths. This avoids any order between branches, at the cost of
an extra $\log$ for sorting.
\end{trick}

\subsection{The centroid tree as a data structure}

The centroid tree also supports \emph{online} queries. Every vertex has $\Oh{\log n}$ centroid ancestors, and by the key idea
above, the path from $v$ to any vertex $w$ passes through one of them. Store at each centroid $c$ a summary of its component
(for example, the distance from $c$ to the nearest marked vertex), store at each vertex the distances to its centroid ancestors,
and a query combines $\Oh{\log n}$ summaries.

\begin{problem}{Nearest red vertex}
In a tree of $n$ vertices, vertex 1 is red and the others are blue. Process $q$ operations: paint vertex $v$ red, or report the
distance from $v$ to the nearest red vertex ($n, q \le 2 \cdot 10^5$). This is Codeforces 342E, \emph{Xenia and Tree}, with
larger limits.
\solution For every centroid $c$ keep $\mathit{best}(c)$, the distance from $c$ to the nearest red vertex of its component. Painting
$v$ updates $\mathit{best}$ at the $\Oh{\log n}$ centroid ancestors of $v$; a query returns $\min_c\, \mathrm{dist}(v, c) +
\mathit{best}(c)$ over them. Every candidate is the length of a real walk, so it is never too small; and the path to the nearest
red vertex $w$ passes through the centroid LCA $c$ of $v$ and $w$, where $\mathrm{dist}(v, c) + \mathit{best}(c) \le
\mathrm{dist}(v, w)$.
\end{problem}

\begin{code}
void solve() {
    int n, q;
    cin >> n >> q;
    vvi adj(n + 1);
    rep(i, n - 1) {
        int a, b;
        cin >> a >> b;
        adj[a].push_back(b);
        adj[b].push_back(a);
    }
    vector<bool> gone(n + 1, false);
    vi par(n + 1), sub(n + 1), dist(n + 1);
    vector<vector<pii>> anc(n + 1);  // (centroid ancestor, distance to it)
    auto bfs = [&](int s, int p) {
        vi order{s};
        par[s] = p;
        rep(i, sz(order)) for (int v : adj[order[i]]) {
            if (v == par[order[i]] || gone[v]) continue;
            par[v] = order[i];
            order.push_back(v);
        }
        return order;
    };
    vi roots{1};
    while (!roots.empty()) {
        int r = roots.back();
        roots.pop_back();
        vi order = bfs(r, 0);
        per(i, sz(order)) {
            int v = order[i];
            sub[v] = 1;
            for (int u : adj[v])
                if (u != par[v] && !gone[u]) sub[v] += sub[u];
        }
        int total = sz(order), c = r;
        for (bool moved = true; moved;) {
            moved = false;
            for (int u : adj[c])
                if (u != par[c] && !gone[u] && 2 * sub[u] > total) {
                    c = u, moved = true;
                    break;
                }
        }
        vi comp = bfs(c, 0);  // c's component, now in BFS order from c
        for (int v : comp) {
            dist[v] = v == c ? 0 : dist[par[v]] + 1;
            anc[v].push_back({c, dist[v]});
        }
        gone[c] = true;
        for (int u : adj[c])
            if (!gone[u]) roots.push_back(u);
    }
    // best[c]: distance from centroid c to the nearest red vertex of its part
    vi best(n + 1, INF);
    auto paint = [&](int v) {
        for (auto [c, d] : anc[v]) chmin(best[c], d);
    };
    paint(1);
    while (q--) {
        int type, v;
        cin >> type >> v;
        if (type == 1) {
            paint(v);
            continue;
        }
        ll res = INF;  // the path to the nearest red vertex passes some c
        for (auto [c, d] : anc[v]) chmin(res, d + best[c]);
        print(res);
    }
}
\end{code}

In the running example, vertex 11 stores its centroid ancestors 0, 4, 7 and 11 at distances 4, 2, 1 and 0. If 9 is red, then
$\mathit{best}(0) = 3$, and the query from 11 finds $4 + 3 = 7$, the true distance, through the centroid 0 where the two vertices
were separated. The lists hold $\Oh{n \log n}$ pairs in total; with $n = q = 2 \cdot 10^5$ the program runs in 0.44\,s on a path
and 0.42\,s on a random tree.

\usetemplate{Centroid decomposition}{Ctrl+Alt+Y, C}{cpcentroid}

\begin{recognize}
``Count or optimise over all paths'' with a condition on length, weight, XOR or colour; ``for each $k$, how many paths of length
$k$''; or online queries about distances to a changing set of vertices. If the answer for a path depends only on the two
endpoints' data relative to a common vertex, centroid decomposition applies.
\end{recognize}

% ======================================================================================================================
\section[DSU on tree]{DSU on tree (Sack) \hfill \plat}\label{sec:sack}

Small-to-large merging (Section~\ref{sec:smalltolarge}) keeps one set per vertex and moves the smaller sets into the larger. When
the information is a global counting array, such as ``how many vertices of each colour'', moving it is impossible and copying it
is too slow. \defn{DSU on tree}, also called \emph{Sack}, keeps a single global array and reuses the heavy child's contents instead
of moving anything. To process a vertex $v$, and either keep or erase its counts afterwards:
\begin{enumerate}[itemsep=1pt]
  \item process each light child, erasing its counts afterwards;
  \item process the heavy child, \emph{keeping} its counts;
  \item add $v$ and every vertex of the light subtrees, walking their Euler-tour ranges;
  \item the array now describes exactly the subtree of $v$: answer $v$'s question;
  \item if $v$ itself is a light child, remove its whole subtree again.
\end{enumerate}

\emph{Why $\Oh{n \log n}$.} A vertex $x$ is added in step 3 once for every ancestor that reaches it through a light child, that
is, once per light edge on its root path, and it is removed in step 5 at most as often. By the light-edge lemma of
Section~\ref{sec:hld} that is at most $\log_2 n$ times, plus once for itself (Figure~\ref{fig:sack}).

\begin{figure}[htbp]
\centering
\begin{tikzpicture}
  \node[vtx, hot] (v) at (0, 0) {$v$};
  \node[vtx, cool] (h) at (0, -1.2) {};
  \node[vtx, warn] (l1) at (-2.1, -1.2) {};
  \node[vtx, warn] (l2) at (2.1, -1.2) {};
  \draw[hed] (v) -- (h);
  \draw[xed] (v) -- (l1) (v) -- (l2);
  \draw[draw=fgblue, fill=fgsky!18, line width=0.6pt] (h.south) -- ++(-0.95, -1.7) -- ++(1.9, 0) -- cycle;
  \draw[draw=fgorange!85!black, fill=fgorange!12, line width=0.6pt] (l1.south) -- ++(-0.5, -0.9) -- ++(1.0, 0) -- cycle;
  \draw[draw=fgorange!85!black, fill=fgorange!12, line width=0.6pt] (l2.south) -- ++(-0.6, -1.1) -- ++(1.2, 0) -- cycle;
  \node[font=\sffamily\scriptsize, text=fgblue, align=center, anchor=north] at (0, -3.33) {heavy child:\\counts kept};
  \node[font=\sffamily\scriptsize, text=fgorange!85!black, align=center, anchor=north] at (-2.1, -2.53) {light child:\\counted,
    erased,\\added again};
  \node[font=\sffamily\scriptsize, text=fgorange!85!black, align=center, anchor=north] at (2.1, -2.73) {light child:\\counted,
    erased,\\added again};
  \begin{scope}[xshift=5.6cm]
    \foreach \nm/\py in {r/0.3, u4/-0.45, u3/-1.2, u2/-1.95, u1/-2.7, x/-3.45}
      \node[vtxs] (\nm) at (0, \py) {};
    \node[font=\sffamily\scriptsize] at (x) {$x$};
    \draw[hed] (r) -- (u4) (u3) -- (u2) (u1) -- (x);
    \draw[xed] (u4) -- (u3) (u2) -- (u1);
    \node[lbl, anchor=west] at (u4.east) {adds $x$ again};
    \node[lbl, anchor=west] at (u2.east) {adds $x$ again};
    \node[notel, anchor=west, text width=3.7cm] at (2.45, -1.55) {Only an ancestor that reaches $x$ through a light child
      (dashed edge) adds $x$ again, so $x$ is added at most $\log_2 n + 1$ times.};
  \end{scope}
\end{tikzpicture}
\caption{DSU on tree. Left: when $v$ is processed, the heavy child's counts are still in the array, and only $v$ and the light
subtrees are added. Right: why every vertex is added $\Oh{\log n}$ times.}
\label{fig:sack}
\end{figure}

\begin{problem}{Dominating colours}
Every vertex of a tree rooted at 1 has a colour in $1..n$. For every vertex $v$, report the sum of all colours that occur in $v$'s
subtree the maximum number of times ($n \le 10^5$; Codeforces 600E, \emph{Lomsat gelral}).
\solution DSU on tree with a counting array $\mathit{cnt}$, the current largest count, and the sum of the colours reaching it.
Adding a vertex updates all three in $\Oh{1}$; erasing resets the two numbers and decrements the counts. The recursion is replaced
by an explicit stack of (vertex, keep, children done) triples, since a path of $10^5$ vertices would otherwise recurse $10^5$
deep.
\end{problem}

\begin{code}
void solve() {
    int n;
    cin >> n;
    vi col(n + 1);
    rep(v, 1, n + 1) cin >> col[v];
    vvi adj(n + 1);
    rep(i, n - 1) {
        int a, b;
        cin >> a >> b;
        adj[a].push_back(b);
        adj[b].push_back(a);
    }
    vi order{1}, par(n + 1, 0), sub(n + 1, 1), heavy(n + 1, 0);
    rep(i, n) for (int v : adj[order[i]])
        if (v != par[order[i]]) par[v] = order[i], order.push_back(v);
    per(i, 1, n) sub[par[order[i]]] += sub[order[i]];
    for (int v : order)
        for (int c : adj[v])
            if (c != par[v] && (!heavy[v] || sub[c] > sub[heavy[v]]))
                heavy[v] = c;
    vi tin(n + 1), ver, st{1};  // preorder: ver[tin[v] .. + sub[v]) = subtree
    while (!st.empty()) {
        int v = st.back();
        st.pop_back();
        tin[v] = sz(ver), ver.push_back(v);
        for (int c : adj[v])
            if (c != par[v]) st.push_back(c);
    }
    vi cnt(n + 1, 0), ans(n + 1);
    ll top = 0, sum = 0;  // the largest count, the sum of colours reaching it
    auto add = [&](int v) {
        int c = col[v];
        if (++cnt[c] > top) top = cnt[c], sum = c;
        else if (cnt[c] == top) sum += c;
    };
    // an explicit stack of (vertex, keep its counts?, children done?)
    vector<array<ll, 3>> todo{{1, 1, 0}};
    while (!todo.empty()) {
        auto [v, keep, done] = todo.back();
        todo.pop_back();
        if (!done) {  // light children run first (popped first), heavy last
            todo.push_back({v, keep, 1});
            if (heavy[v]) todo.push_back({heavy[v], 1, 0});
            for (int c : adj[v])
                if (c != par[v] && c != heavy[v]) todo.push_back({c, 0, 0});
            continue;
        }
        // the heavy child's counts are still there; add everything else
        for (int c : adj[v])
            if (c != par[v] && c != heavy[v])
                rep(t, tin[c], tin[c] + sub[c]) add(ver[t]);
        add(v);
        ans[v] = sum;
        if (!keep) {  // a light child: leave the counts empty again
            rep(t, tin[v], tin[v] + sub[v]) cnt[col[ver[t]]]--;
            top = sum = 0;
        }
    }
    print(vi(ans.begin() + 1, ans.end()));
}
\end{code}

\usetemplate{DSU on tree (Sack)}{Ctrl+Alt+U, Shift+T}{cpsack}

The order of the stack matters: the heavy child is pushed \emph{before} the light children, so it is popped, and processed,
after all of them, and its counts are the ones still present when $v$ is finished. With $n = 2 \cdot 10^5$ the program runs in
0.17\,s on a path and 0.12\,s on a random tree.

\begin{tip}
Small-to-large merging, DSU on tree and the offline methods of later chapters overlap. Use small-to-large when each vertex needs a
set or map of its own afterwards; DSU on tree when the statistic lives in one array (counts by colour, by depth, by value) and
must support ``the most frequent'' or ``how many with count at least $k$''; and the Euler tour with an offline range technique when
the question is a plain range question on the flattened tree, such as the number of distinct values (Mo's algorithm,
Chapter~\ref{ch:advds}).
\end{tip}

% ======================================================================================================================
\section[Long-path decomposition]{Long-path decomposition \hfill \plat}\label{sec:longpath}

Many tree DPs are indexed by depth: $f_v[d]$ is the number (or best value) of vertices at distance $d$ below $v$. Merging the
children's arrays naively costs the sum of all heights over all vertices, which is $\Theta(n^2)$ on a path. Small to large by
height gives $\Oh{n \log n}$; a sharper idea gives $\Oh{n}$.

\paragraph{Long paths.} The \defn{height} $h(v)$ is the number of vertices on the longest downward path from $v$. Call the child
with the largest height the \emph{long child} and decompose the tree into \defn{long paths} exactly as HLD does with sizes. Now
$f_v[d] = f_{\mathit{long}}[d - 1]$ for every $d \ge 1$, plus the contributions of the short children. So give each long path one
block of a shared buffer, as long as the path, and let the array of every vertex on the path start one cell after its parent's
(Figure~\ref{fig:longpath}): a vertex inherits its long child's array without copying, sets $f_v[0] = 1$, and merges only its short
children by hand.

\emph{Why $\Oh{n}$.} Merging a short child $c$ costs $h(c)$, the length of the long path that starts at $c$. Every long path
starts at a short child (or at the root) and is merged once, into its parent, so the total cost is the total length of all long
paths, which is $n$.

\begin{figure}[htbp]
\centering
\begin{tikzpicture}
  \foreach \nm/\px/\py/\sy in {1/1.8/0/hot, 2/1.0/-1.0/hot, 3/2.6/-1.0/cool, 4/0.5/-2.0/hot, 5/1.5/-2.0/n, 8/2.6/-2.0/cool,
      6/0.5/-3.0/hot, 9/2.6/-3.0/cool, 7/0.5/-4.0/hot}
    \node[vtx, \sy] (L\nm) at (\px, \py) {\nm};
  \draw[hed] (L1) -- (L2) -- (L4) -- (L6) -- (L7);
  \draw[bed] (L3) -- (L8) -- (L9);
  \draw[xed] (L1) -- (L3) (L2) -- (L5);
  \foreach \ci/\nm in {1/1, 2/2, 3/4, 4/6, 5/7, 6/3, 7/8, 8/9, 9/5}
    \node[font=\sffamily\scriptsize\bfseries, text=fgink] at ({4.2 + 0.7*\ci - 0.35}, -0.3) {\nm};
  \node[lbl, anchor=east] at (4.05, -0.3) {$f[0]$ of};
  \node[blbl, anchor=east] at (4.05, -1.0) {buffer};
  \arrs{B}{4.2}{-1.0}{1/hot, 2/hot, 3/hot, 2/hot, 1/hot, 1/cool, 1/cool, 1/cool, 1/n}
  \foreach \nm/\fc/\yy in {1/1/-1.62, 2/2/-1.92, 4/3/-2.22} {
    \draw[line width=1.8pt, draw=fgverm!75, line cap=butt] ({4.25 + 0.7*(\fc - 1)}, \yy) -- (7.65, \yy);
    \node[font=\sffamily\scriptsize, text=fgverm, anchor=east] at ({4.2 + 0.7*(\fc - 1)}, \yy) {$f_{\nm}$};
  }
  \draw[line width=1.8pt, draw=fgblue!75, line cap=butt] (7.75, -1.62) -- (9.75, -1.62);
  \node[font=\sffamily\scriptsize, text=fgblue, anchor=west] at (9.8, -1.62) {$f_3$};
  \draw[->, line width=1pt, draw=fgblue] ($(B-6.north)+(0, 0.62)$) to[bend right=35]
    node[above, font=\sffamily\scriptsize, text=fgblue] {merge 3: $f_1[d + 1] \mathrel{+}= f_3[d]$} ($(B-2.north)+(0, 0.62)$);
  \node[notel, anchor=north west, text width=6.0cm] at (4.2, -2.6) {Long paths 1--2--4--6--7 (red), 3--8--9 (blue) and 5. The
    bars show where the arrays of 1, 2 and 4 live: each starts one cell after its parent's. The final block of 1 reads
    $f_1 = (1, 2, 3, 2, 1)$.};
\end{tikzpicture}
\caption{Long-path decomposition. Each long path owns one block of the buffer, so a vertex inherits its long child's counts in
$\Oh{1}$; only the short child 3 is merged by hand. Three vertices lie at distance 2 below 1, the most at any distance, so the
dominant depth of 1 is 2.}
\label{fig:longpath}
\end{figure}

\begin{problem}{Dominant depths}
In a tree rooted at 1 with $n \le 10^6$ vertices, for every vertex $v$ find the smallest $d$ that maximises the number of
vertices of $v$'s subtree at distance exactly $d$ from $v$ (Codeforces 1009F, \emph{Dominant Indices}).
\solution Long-path decomposition with counts. Vertex $v$ inherits the best depth of its long child plus one, compares it with
$f_v[0] = 1$, and re-checks only the cells that its short children change; values only grow, so the running best stays correct.
\end{problem}

\begin{code}
void solve() {
    int n;
    cin >> n;
    vvi adj(n + 1);
    rep(i, n - 1) {
        int a, b;
        cin >> a >> b;
        adj[a].push_back(b);
        adj[b].push_back(a);
    }
    vi order{1}, par(n + 1, 0);
    rep(i, n) for (int v : adj[order[i]])
        if (v != par[order[i]]) par[v] = order[i], order.push_back(v);
    // h[v]: vertices on the longest downward path from v; lc[v]: its child
    vi h(n + 1, 1), lc(n + 1, 0);
    per(i, 1, n) {
        int v = order[i], p = par[v];
        if (h[v] + 1 > h[p]) h[p] = h[v] + 1, lc[p] = v;
    }
    // each long path owns one block of the buffer, and along the path the
    // arrays are shifted by one: a vertex inherits its long child's array
    vi pos(n + 1), cnt(n, 0), ans(n + 1);
    int next = 0;
    for (int v : order) {
        if (v == 1 || lc[par[v]] != v) pos[v] = next, next += h[v];
        else pos[v] = pos[par[v]] + 1;
    }
    per(i, n) {  // children first
        int v = order[i];
        auto f = cnt.begin() + pos[v];  // f[d]: vertices at distance d below v
        f[0] = 1;
        ans[v] = lc[v] ? ans[lc[v]] + 1 : 0;
        if (f[ans[v]] <= 1) ans[v] = 0;
        for (int c : adj[v]) {  // merge the short children by hand
            if (c == par[v] || c == lc[v]) continue;
            rep(d, h[c]) {
                f[d + 1] += cnt[pos[c] + d];
                ll a = f[d + 1], b = f[ans[v]];
                if (a > b || (a == b && d + 1 < ans[v])) ans[v] = d + 1;
            }
        }
    }
    print(vi(ans.begin() + 1, ans.end()));
}
\end{code}

\usetemplate{Long-path decomposition}{Ctrl+Alt+U, Shift+N}{cplongpath}

The buffer has exactly $n$ cells, because the blocks are as long as the long paths and those partition the vertices. With $n =
10^6$ the program runs in 0.42\,s on a path, 0.39\,s on a random tree and 0.36\,s on a caterpillar, most of it spent reading the
input.

\begin{trick}[Level ancestors in $\Oh{1}$: the ladder algorithm]
Extend every long path upwards by its own length to get a \emph{ladder} of at most twice the length; the ladders have $\Oh{n}$
cells in total. To find the $k$-th ancestor of $v$, first jump $2^j$ levels with binary lifting, where $2^j$ is the highest power
of two not above $k$, landing at $u$. Since $u$ has a descendant $2^j$ levels below it, the long path through $u$ continues for at
least $2^j$ more vertices below $u$, so its ladder extends at least $2^j > k - 2^j$ levels above its head, and the remaining
climb is a single array lookup in that ladder. Preprocessing is $\Oh{n \log n}$, and every query is two lookups.
\end{trick}

% ======================================================================================================================
\section[Virtual trees]{Virtual (auxiliary) trees \hfill \plat}\label{sec:virtualtree}

Some problems ask many queries, each about a small set of $k$ \emph{key vertices}, with the sum of all $k$ bounded by about $2
\cdot 10^5$. A DP over the whole tree costs $\Oh{n}$ per query, far too much, but most of the tree is irrelevant to a query: only
the keys and the places where their paths branch matter. The \defn{virtual tree} of the keys contains the keys and the LCAs of
pairs of keys; its edges stand for the compressed paths between them, with the path length as weight
(Figure~\ref{fig:virtualtree}).

\begin{figure}[htbp]
\centering
\begin{tikzpicture}
  \rexcoords
  \draw[line width=5pt, draw=fggreen, opacity=0.3] (R11) -- (R7) -- (R4) -- (R1) -- (R0) -- (R2) -- (R6) -- (R10) -- (R12)
    (R4) -- (R8) (R1) -- (R5);
  \draw[ed] \rexedges;
  \foreach \nm/\sy in {0/cool, 1/cool, 2/n, 3/n, 4/cool, 5/hot, 6/n, 7/n, 8/hot, 9/n, 10/n, 11/hot, 12/hot}
    \node[vtx, \sy] at (R\nm) {\nm};
  \begin{scope}[xshift=8.6cm]
    \foreach \nm/\px/\py/\sy in {0/1.6/0/cool, 1/0.6/-1.3/cool, 12/2.6/-1.3/hot, 4/0.0/-2.6/cool, 5/1.2/-2.6/hot,
        11/-0.5/-3.9/hot, 8/0.5/-3.9/hot}
      \node[vtx, \sy] (W\nm) at (\px, \py) {\nm};
    \foreach \pa/\ch/\elen in {0/1/1, 0/12/4, 1/4/1, 1/5/1, 4/11/2, 4/8/1}
      \draw[ed] (W\pa) -- node[wl, font=\sffamily\scriptsize\bfseries, text=fggreen!50!black] {\elen} (W\ch);
    \node[lbl] at (1.0, -4.75) {virtual tree; edge weights are path lengths};
  \end{scope}
\end{tikzpicture}
\caption{The virtual tree of the keys 11, 8, 5 and 12 (red). In DFS order the keys are 11, 8, 5, 12, and the LCAs of neighbours
in that order, 4, 1 and 0 (blue), complete the vertex set. The smallest subtree containing the keys (green) has 10 edges, the
total weight of the virtual tree.}
\label{fig:virtualtree}
\end{figure}

\paragraph{Construction.} Sort the keys by DFS entry time, $v_1, \dots, v_k$. Every LCA of two keys is already the LCA of two
\emph{neighbours} in this order: for $i < j$, $\mathrm{lca}(v_i, v_j)$ is the shallowest of $\mathrm{lca}(v_t, v_{t+1})$ for $i \le t
< j$. So adding the $k - 1$ neighbour LCAs closes the set under LCA, and the virtual tree has at most $2k - 1$ vertices. Sort this
set by entry time again, $w_1, \dots, w_m$; then the virtual parent of $w_i$ is $\mathrm{lca}(w_{i-1}, w_i)$.

\emph{Why the parent rule works.} The parent of $w_i$ must be the deepest vertex of the set that is a proper ancestor of $w_i$.
The vertex $\mathrm{lca}(w_{i-1}, w_i)$ belongs to the set and is an ancestor of $w_i$; it is not $w_i$ itself, since $w_i$ comes
after $w_{i-1}$ in preorder and so cannot be its ancestor. Any other ancestor $a$ of $w_i$ in the set is entered before $w_i$,
hence no later than $w_{i-1}$, and its subtree range covers everything up to $w_i$, so $a$ is an ancestor of $w_{i-1}$ too, and
therefore of their LCA. With $\Oh{\log n}$ LCAs, one query costs $\Oh{k \log n}$ after the sorting.

\begin{problem}{Key vertex queries}
A tree has $n$ vertices. Each of $q$ queries gives $k$ distinct key vertices ($\sum k \le 2 \cdot 10^5$); report the number of
edges of the smallest subtree containing all keys, and the sum of the distances over all pairs of keys.
\solution Build the virtual tree; the HLD template provides both a DFS order (\cc{pos}) and LCAs. Walk the virtual vertices in
reverse DFS order, children first, and count the keys below each one. A virtual edge of length $\ell$ above a part with $b$ keys
belongs to the smallest subtree, and every one of its $\ell$ original edges separates $b$ keys from the other $k - b$, so it adds
$\ell$ to the first answer and $\ell\, b\, (k - b)$ to the second (the ``count what each edge separates'' trick of
Section~\ref{sec:rerooting}).
\end{problem}

\begin{code}
void solve() {  // cphld and cpvirtual are pasted above solve()
    int n, q;
    cin >> n >> q;
    vvi adj(n);
    rep(i, n - 1) {
        int a, b;
        cin >> a >> b;
        a--, b--;
        adj[a].push_back(b);
        adj[b].push_back(a);
    }
    HLD h(adj);  // h.pos is a DFS order, h.lca an LCA
    auto lca = [&](int a, int b) { return h.lca(a, b); };
    vi idx(n, -1), below;
    while (q--) {
        int k;
        cin >> k;
        vi keys(k);
        for (ll& v : keys) cin >> v, v--;
        auto [vs, vp] = virtual_tree(keys, h.pos, lca);
        int m = sz(vs);
        rep(i, m) idx[vs[i]] = i;
        below.assign(m, 0);  // key vertices in each virtual subtree
        for (int v : keys) below[idx[v]] = 1;
        ll steiner = 0, pairs = 0;
        per(i, 1, m) {  // vs is in DFS order: children before parents
            int v = vs[i], p = vp[i];
            ll len = h.depth[v] - h.depth[p];  // edges of the compressed path
            steiner += len;
            pairs += len * below[i] * (k - below[i]);
            below[idx[p]] += below[i];
        }
        print(steiner, pairs);
        for (int v : vs) idx[v] = -1;
    }
}
\end{code}

\usetemplate{Virtual tree}{Ctrl+Alt+Y, V}{cpvirtual}

The array \cc{idx} maps a vertex to its index in the current virtual tree and is reset after each query, touching only the
$\Oh{k}$ vertices used; clearing an $n$-sized array per query would cost $\Oh{nq}$, the classic mistake with virtual trees. With
$\sum k = 2 \cdot 10^5$ the program runs in 0.15\,s.

\begin{trick}[The smallest subtree from the DFS order alone]
For keys sorted by entry time, $v_1, \dots, v_k$, the number of edges of the smallest subtree containing them is
\[
  \tfrac12 \bigl(\mathrm{dist}(v_1, v_2) + \mathrm{dist}(v_2, v_3) + \dots + \mathrm{dist}(v_{k-1}, v_k) + \mathrm{dist}(v_k,
  v_1)\bigr),
\]
because walking from key to key in DFS order and back to the start traverses every edge of that subtree exactly twice, like the
DFS itself. In Figure~\ref{fig:virtualtree}: $(3 + 3 + 6 + 8)/2 = 10$. With a set ordered by entry time, the formula also maintains
the subtree size under insertions and deletions of keys in $\Oh{\log n}$ each.
\end{trick}

\begin{recognize}
Many queries, each naming a few vertices, with ``the sum of $k$ over all queries is at most \dots'' in the constraints, and a
question that would be an easy tree DP if it were asked once.
\end{recognize}

% ======================================================================================================================
\section[Euler tour trees]{Euler tour trees: link and cut \hfill \ioi}\label{sec:ett}

So far the tree never changed. In a \emph{dynamic forest}, edges are added (link) and removed (cut) online, and queries ask
whether two vertices are connected, or for the sum of the values in a vertex's tree. When all operations are known in advance,
divide and conquer over time with a rollback DSU is usually simpler (Section~\ref{sec:offline-dyncon}); online, keep every tree as
its Euler tour inside a balanced binary search tree.

\paragraph{Representation.} An \defn{Euler tour tree} stores the tour of each tree as a sequence that contains every vertex once
and every edge twice, once in each direction (Figure~\ref{fig:ett}). The sequence is read cyclically, so rerooting the tree at $v$
is a rotation that starts the sequence at $v$: one split and one merge. With a treap (Chapter~\ref{ch:advds}) every operation is
$\Oh{\log n}$ expected:
\begin{itemize}[itemsep=1pt]
  \item \textbf{connected$(u, v)$}: are $u$ and $v$ in the same treap? Compare the roots.
  \item \textbf{link$(u, v)$}: reroot both trees at $u$ and $v$, then concatenate tour$(u)$, the node $(u, v)$, tour$(v)$ and the
        node $(v, u)$.
  \item \textbf{cut$(u, v)$}: reroot at $u$; the tour is now $A\,(u, v)\,B\,(v, u)\,C$, where $B$ is $v$'s side. Split out $B$
        and the two edge nodes, and merge $A$ with $C$.
  \item \textbf{aggregates}: every treap node stores the sum (or minimum, or count of vertices) of its subtree, so the treap root
        holds the aggregate of the whole tree; changing a vertex value updates one root path.
\end{itemize}

\begin{figure}[htbp]
\centering
\begin{tikzpicture}
  \node[vtx] (ta) at (0, 0) {a}; \node[vtx] (tb) at (0, -1.0) {b};
  \node[vtx] (tc) at (-0.6, -2.0) {c}; \node[vtx, cool] (td) at (0.6, -2.0) {d};
  \draw[ed] (ta) -- (tb) (tb) -- (tc);
  \draw[hed] (tb) -- (td);
  \node[blbl, anchor=east] at (4.3, 0) {tour from a};
  \arrs{T}{4.45}{0}{a/n, ab/dim, b/n, bc/dim, c/n, cb/dim, bd/dim, d/n, db/dim, ba/dim}
  \node[blbl, anchor=east] at (4.3, -1.0) {rerooted at b};
  \arrs{U}{4.45}{-1.0}{b/n, bc/dim, c/n, cb/dim, bd/hot, d/cool, db/hot, ba/dim, a/n, ab/dim}
  \cellframe{U-5}{U-7}{fgverm}
  \node[blbl, anchor=east] at (4.3, -2.0) {after cutting b--d};
  \arrs{X}{4.45}{-2.0}{b/n, bc/dim, c/n, cb/dim, ba/dim, a/n, ab/dim}
  \arrs{Y}{10.05}{-2.0}{d/cool}
\end{tikzpicture}
\caption{An Euler tour tree. The tour lists every vertex once (white) and every edge once in each direction (grey). Rerooting at b
rotates the tour; cutting b--d removes the two edge nodes (red) together with d's side between them (blue), which leaves the tour
of the remaining tree and the tour of d alone.}
\label{fig:ett}
\end{figure}

\begin{problem}{Forest operations}
A forest on $n$ vertices starts without edges; every vertex has a value. Process online: link two vertices of different trees,
cut an existing edge, ask whether two vertices are connected, change a value, and report the sum of the values in a vertex's tree.
\end{problem}

\begin{code}
void solve() {  // Euler tour tree (cpett) is pasted above solve()
    int n, q;
    cin >> n >> q;
    ETT f(n);
    rep(v, n) {
        ll x;
        cin >> x;
        f.set(v, x);
    }
    while (q--) {
        int type, u;
        cin >> type >> u;
        u--;
        if (type == 5) {  // the sum of the values in u's tree
            print(f.comp_sum(u));
            continue;
        }
        ll x;
        cin >> x;
        if (type == 1) f.link(u, x - 1);
        else if (type == 2) f.cut(u, x - 1);
        else if (type == 3) yn(f.connected(u, x - 1));
        else f.set(u, x);
    }
}
\end{code}

\usetemplate{Euler tour tree}{Ctrl+Alt+U, E}{cpett}

The template finds the treap node of an edge through a \cc{map} from vertex pairs, which costs another $\Oh{\log n}$ per link or
cut. Subtree aggregates are available too: to get the subtree of $v$ when the tree is rooted at $r$, reroot at $r$, and the subtree
is the part of the tour between the nodes $(p, v)$ and $(v, p)$, where $p$ is the parent of $v$.

% ======================================================================================================================
\section[Link-cut trees and top trees]{Link-cut trees and top trees \hfill \ioi}\label{sec:lct}

Euler tour trees answer questions about whole trees and subtrees, but not about \emph{paths}: the tour mixes the vertices of a path
with everything hanging from it. A \defn{link-cut tree} (Sleator and Tarjan) keeps a dynamic forest with path aggregates, such as
the maximum edge weight between two vertices, in $\Oh{\log n}$ amortised time per operation. It splits each tree into
\emph{preferred paths}, much like the chains of HLD but changing over time, and stores each preferred path in a splay tree ordered
by depth. Its central operation, \emph{access}$(v)$, makes the path from the root to $v$ preferred, after which that path is a
single splay tree and any path aggregate is read at its root. Rerooting reverses a whole preferred path, with a lazy ``reverse''
flag.

The classic use is a \emph{dynamic minimum spanning forest} under edge insertions: when an edge $u$--$v$ of weight $w$ arrives and
$u, v$ are already connected, find the heaviest edge on the path between them; if it is heavier than $w$, cut it and link the new
edge, which by the cycle property of Section~\ref{sec:mst} keeps the forest minimum. \defn{Top trees} combine path and subtree
aggregates in one structure, with $\Oh{\log n}$ per operation, but they are long to implement and rarely appear outside the
hardest olympiad problems. Table~\ref{tab:dyntrees} summarises the structures of this chapter.

\begin{table}[htbp]
\centering\small
\begin{tabular}{@{}lccccl@{}}
\toprule
structure & link, cut & path aggregates & subtree aggregates & time per operation & code \\
\midrule
HLD + segment tree & no & yes & yes & $\Oh{\log^2 n}$ & short \\
centroid tree & no & distances & no & $\Oh{\log n}$ & short \\
Euler tour tree & yes & no & yes & $\Oh{\log n}$ expected & medium \\
link-cut tree & yes & yes & extra work & $\Oh{\log n}$ amortised & medium \\
top tree & yes & yes & yes & $\Oh{\log n}$ & long \\
\bottomrule
\end{tabular}
\caption{What the tree structures of this chapter support. ``Distances'' means queries that combine distances to centroid
ancestors, as in the nearest-red-vertex problem.}
\label{tab:dyntrees}
\end{table}

\begin{tip}
Before reaching for a dynamic tree, check whether the problem is offline. Operations known in advance usually yield to a sweep, to
divide and conquer over time with a rollback DSU, or to a static HLD over the final tree, all of them shorter and safer to write
in a contest.
\end{tip}

\begin{summary}[Chapter summary]
\begin{itemize}
  \item A root path crosses at most $\log_2 n$ light edges, so HLD cuts every path into $\Oh{\log n}$ chain pieces, each a range of
        positions; positions also keep subtrees contiguous.
  \item Centroid decomposition has $\Oh{\log n}$ levels; every path is handled at exactly one centroid, its LCA in the centroid
        tree, which gives $\Oh{n \log n}$ algorithms for all-paths problems and $\Oh{\log n}$ online distance queries.
  \item DSU on tree keeps the heavy child's counts and re-adds light subtrees: $\Oh{n \log n}$ with a single global array.
  \item Long-path decomposition shares one buffer along each long path and runs depth-indexed DP in $\Oh{n}$.
  \item A virtual tree on $k$ keys has at most $2k - 1$ vertices: sort by entry time, add neighbour LCAs, and the parent of each
        vertex is its LCA with its predecessor.
  \item Euler tour trees keep a forest under link and cut with component aggregates; link-cut trees add path aggregates.
\end{itemize}
\end{summary}

\begin{practice}
\ppgroup{Heavy-light and centroid decomposition}
\begin{pplist}
  \item Path Queries II
  \item Counting Paths
  \item Finding a Centroid
  \item Fixed-Length Paths I
  \item Fixed-Length Paths II
\end{pplist}
\ppgroup{Subtree statistics}
\begin{pplist}
  \item Distinct Colors
  \item Subtree Queries
  \item Path Queries
\end{pplist}
\end{practice}
